\section{CTI Workflows}

Cyber Threat Intelligence becomes most valuable when it directly supports operational teams. Effective CTI workflows bridge the gap between intelligence reporting and hands-on defensive work by transforming threat insights into practical actions. For a CTI analyst, threat hunter, or detection engineer, the goal is not simply to collect intelligence but to operationalize it—turning adversary behaviors into hypotheses, detections, dashboards, and investigative leads. This section outlines the core workflows that enable CTI to drive meaningful security outcomes.

A typical CTI workflow begins with the identification of relevant intelligence. This may come from threat reports, malware analysis, internal telemetry, or external feeds. The analyst evaluates the intelligence to determine what behaviors, techniques, or infrastructure are most relevant to the organization. Instead of focusing on short-lived indicators such as IP addresses or domains, the emphasis is placed on techniques—how the adversary gains access, executes code, moves laterally, or evades detection. These behaviors form the foundation for operational work.

Once relevant behaviors are identified, the next step is translation. CTI analysts convert high-level intelligence into concrete questions that can be answered with telemetry. For example, if a threat actor is known to use encoded PowerShell commands, the operational question becomes: “Where in our environment do we see PowerShell with encoded commands?” This translation step is critical because it connects intelligence to the data sources available in Splunk or Sentinel.

The translated intelligence then becomes the basis for threat hunting. Analysts use SPL or KQL to search for behaviors associated with the threat actor. The following example illustrates how CTI-driven intelligence about encoded PowerShell commands becomes a hunt:

\textbf{SPL Example}
\begin{lstlisting}
index=sysmon EventCode=1 Image="*powershell.exe"
| search CommandLine="*EncodedCommand*"
| table _time, Computer, User, ParentImage, Image, CommandLine
\end{lstlisting}

\textbf{KQL Equivalent}
\begin{lstlisting}
Sysmon
| where EventID == 1 and Image endswith "powershell.exe"
| where CommandLine contains "EncodedCommand"
| project TimeGenerated, Computer, User, ParentImage, Image, CommandLine
\end{lstlisting}

If the hunt reveals suspicious activity, the workflow moves into detection engineering. The same behavioral pattern is refined into a repeatable detection that continuously monitors for similar activity. This ensures that the organization is protected not only against the current threat but also against future variants that use the same technique.

CTI workflows also support dashboard creation. Intelligence about adversary behavior informs which panels should be included in dashboards and how they should be structured. For example, if a threat actor frequently uses scheduled tasks for persistence, a dashboard panel that highlights scheduled task creation becomes valuable. CTI ensures that dashboards remain aligned with real-world threats rather than generic system activity.

Another key component of CTI workflows is enrichment. When analysts encounter suspicious events in Splunk or Sentinel, CTI provides context that helps determine whether the activity is benign or malicious. For example, a suspicious domain discovered in DNS logs can be enriched with intelligence about whether it is associated with a known threat actor. This enrichment accelerates investigations and improves decision-making.

Finally, CTI workflows rely on feedback loops. Threat hunters, SOC analysts, and detection engineers provide insight into which intelligence is useful and which detections generate noise. This feedback helps refine intelligence requirements and ensures that CTI remains aligned with operational needs. Over time, this iterative process strengthens the organization’s defensive posture and improves the quality of both intelligence and detections.

The workflows described in this section form the backbone of operational CTI. By connecting intelligence to telemetry, hunts, detections, dashboards, and investigations, CTI becomes a practical force multiplier rather than a theoretical exercise. These workflows will be expanded in later sections into full playbooks and detection patterns tailored to specific adversary techniques.
